\scanpage{81}
\mspage{80}
of $\handwritten{\underline}{Q}_p$ \handwritten{)} which are therefore isomorphic \inserted{\handwritten{``geometrically'' i.e. in $\operatorname{Rep}(\underline{H})$}}\handwritten{\begin{tikzpicture}[baseline=0cm,x=1cm,y=1cm]\draw[thin] (0,-.42)--(0,0)..controls (0,.12) and (.08,.12)..(.18,.12)--(3.1,.17)--(3.1,.28);\end{tikzpicture}} (but probably not canonically) to the $E_{r,1}$ for $r\geq1$: they give canonic\rlap{d}al free modules of rank 1 over the $K_{r,1}$ ...\handwritten{]}

9. For arbitrary $k$, one finds a \deleted{ho} canonical $\handwritten{\otimes}$-homomorphism \deleted{(with the help} defined up to unique isomorphism (one uses a choice of an algebraic closure $\handwritten{\overline}{k}$ of $k$, but this choice is inessential ..)
\[
\operatorname{Fcriso}(k)\handwritten{\longrightarrow}\operatorname{Rep}(\handwritten{\underline}{H})
\]
or again a homomorphism of gerbes
\[
\handwritten{\underline}{H}\handwritten{\longrightarrow}\handwritten{\underline}{G}(k)
\]
(\unsure{\rlap{a}r|r}ecall that this homomorphism factors through the kernel,\deleted{of} namely \deleted{\unsure{Gl|G|Gr}} $\handwritten{\underline}{H}(k)$, of $\handwritten{\underline}{G}(k)\handwritten{\to}\handwritten{\underline{\Pi}}_p(k)$, and that f\rlap{i}or perfect $k$ one guesses that the resulting homomorphism $\handwritten{\underline}{H}\handwritten{\to}\handwritten{\underline}{H}(k)$ is an equivalence ...). If $\handwritten{\lambda\in}\handwritten{\underline}{Q}$, an F\deleted{(}\handwritten{-}isocrystal over $k$ is said to be \deleted{homogen} \handwritten{\uline}{of slope} $\handwritten{\lambda}$ if its image in $\operatorname{Rep}(\handwritten{\underline}{H})$ is of slope $\handwritten{\lambda}$. One then sees at once that an F-isocrystal is of slope zero if and only if it is ``trivial'' i.e. satisfies the conditions of 5.2 (which justifies the terminology which is introduced, a little prematurely, at that point), it is effective if and only if its image in \inserted{$\operatorname{Rep}(H)$ is so,i.e. is of slope $\handwritten{\geq}$ 0.}

Let us consider the associated homomorphism of bands
\[
H\handwritten{\longrightarrow}G(k)\quad ,
\]
one fin\rlap{e}ds th\unsure{\rlap{r}e|e}\unsure{\rlap{n}n|n|n} that \rlap{e}if $k$ is perfect, the preceding homomorphism is \handwritten{\uline}{central}, in other words $H$ acts on the identity functor of $\operatorname{Fcriso}(k)\handwritten{\to}\operatorname{Rep}(\handwritten{\underline}{G}(k))$, so that every F-isocrystal \inserted{$(E,F_E)$}\handwritten{\begin{tikzpicture}[baseline=0cm,x=1cm,y=1cm]\draw[thin] (0,-.18)..controls (.04,.15) and (.38,.12)..(.65,.08)..controls (.9,.05) and (1.25,.16)..(1.3,.24)..controls (1.33,.27) and (1.3,.3)..(1.28,.32);\end{tikzpicture}} decomposes into its \handwritten{``}\handwritten{\uline}{isoclinic components}\handwritten{''} $E_{\handwritten{\lambda}}$, $\handwritten{\lambda\in}\handwritten{\underline}{Q}$. $E$ is ``trivial'' resp. effective if and only if $E_{\handwritten{\lambda}}=0$ for $\handwritten{\lambda\ne}$ 0 (resp. $\handwritten{\lambda<0}$).

Note that the Tate \inserted{F-iso}\handwritten{\begin{tikzpicture}[baseline=0cm,x=1cm,y=1cm]\draw[thin] (0,-.2)..controls (0,.13) and (.38,.14)..(.62,.16)..controls (.89,.18) and (1.02,.33)..(.97,.36);\end{tikzpicture}}crystal $K(\unsure{1|l})$ is of slope\handwritten{-}1. Thus (for perfect $k$) every F\deleted{(}\handwritten{-}isocrystal can be written uniquely in the form
\[
E=\handwritten{\sum_{i\in\mathbb{Z}}}E_i(-i)\quad ,
\]
where $E_i$ has isoclinic components of slope $0\handwritten{<}\,\inserted{\handwritten{\rlap{\raisebox{.5em}{\begin{tikzpicture}[x=1cm,y=1cm]\draw[gray,thin] (-.15,.1)..controls (0,.38) and (.45,.3)..(.55,-.03);\end{tikzpicture}}}i}}\handwritten{\leq}1$. Moreover,
